\chapter{Theory and Analytical Design}

\section{Governing model}
Introduce only the theory needed to justify engineering choices. For antenna and array projects, canonical references such as Stutzman and Thiele provide a disciplined starting point for notation and first-order design relationships~\cite{stutzman2012}. Define variables,
assumptions, coordinate systems, and validity limits. For an illustrative normalized
array or sensing problem, a generic response may be written as
\begin{equation}
H(\theta)=\sum_{n=0}^{N-1} w_n e^{j n \psi(\theta)},
\label{eq:generic-response}
\end{equation}
where $w_n$ represents a complex weighting coefficient and $\psi(\theta)$ represents
the progressive phase associated with geometry and excitation.

\section{First-order design calculation}
Show the calculation chain that converts requirements into initial dimensions or
parameters. The purpose is not to reproduce a textbook chapter; it is to establish a
traceable starting point for modeling and optimization.

\begin{analyticalevidencebox}
The example script in Appendix~\ref{app:analytical-verification} calculates a normalized
wavelength, nominal spacing, and progressive phase. Replace it with the project's real
verification script and retain its console output under \texttt{verification\_records/}.
\end{analyticalevidencebox}

\section{Assumptions and validity}
Record assumptions before presenting results. Typical examples include linearity,
steady-state operation, homogeneous materials, ideal source behavior, negligible
connector effects, limited temperature range, or simplified boundary conditions.
